% ===================================================================
%  TAULA N.º 8 — La Collita. — La Caça, la Verema, la Pesca.
%  Traduction 2026 d'apres texto/io, controlee sur texto/fr et
%  texto/es. Consigne : voir texto/ca/10-taula-01.tex.
%
%  L'appel de note est dans le TITRE de la scene. LE RENVOI (13) DE
%  LA FAUX MANQUE A L'IDO ; la colonne le rend : ecart declare.
% ===================================================================

%%K t08-noto-1 noto
\VUnotes{143.2mm}{%
\textit{(*) Perquè aquest mot no és precís: es tracta d'una collita de lli,
d'una verema, d'una recol·lecció de pomes? Potser fóra millor usar «}
\VUgras{mesono} \textit{», proposat a} Mondo \textit{XIV, p. 242. Però fins ara
aquella arrel nova no ha estat adoptada per l'Acadèmia i espera la discussió
segons les regles habituals del moviment idista.}%
}

%%K t08-apar-1 apar
\VUsaut{5.0mm}
\VUtitre{15.0pt}{60}{TAULA N.º 8}
\VUsaut{3.0mm}
\VUcentre{13.2pt}{90}{\textit{Primera escena.}}
\VUsaut{3.0mm}
\VUpk{10.2pt}{40}{La Collita (*).}
\VUsaut{4.0mm}
\VUpk{10.2pt}{40}{Els Aspectes del Camp.}
\VUsaut{3.0mm}

%%K t08-c1-01-1 p
1. --- Amb l'\VUgras{estiu}, el camp ha canviat d'aspecte una vegada més. La
llavor confiada al solc ha germinat; una planteta verda ha sortit de la terra i
ha donat la seva espiga. El gra s'ha format, després ha madurat, i ja no queda
més que collir.

%%K t08-c1-02-1 p
2. --- Les \VUgras{flors del camp} són obertes. \VUgras{Blauets}
\textsuperscript{(1)}, \VUgras{roselles} \textsuperscript{(2)} i
\VUgras{didaleres} \textsuperscript{(3)} posen en el to uniforme dels blats
l'alegre contrast de les seves fresques \VUgras{corol·les}.

%%K t08-c1-03-1 p
3. --- Els arbres fruiters s'han cobert de fulles; el fruit ha madurat. El petit
\VUgras{cirerer} \textsuperscript{(4)} és carregat de belles \VUgras{cireres}
\textsuperscript{(5)}, que els nens cullen i mengen amb delit.

%%K t08-c1-04-1 p
4. --- El ramat ja pot passar tot el dia a l'aire lliure.

%%K t08-c1-04-2 p
Els \VUgras{anyells} \textsuperscript{(6)} han crescut i s'han convertit en belles
\VUgras{ovelles} \textsuperscript{(7)} i bells \VUgras{marrans}
\textsuperscript{(8)} d'espès velló. Pasturen tranquil·lament l'\VUgras{herba} de
la \VUgras{pastura} \textsuperscript{(9)}, esquitxada de \VUgras{margarides}.

%%K t08-c1-04-3 p
Aviat començaran a tondre aquests animals per tenir-ne la \VUgras{llana}, amb la
qual es fabrica el drap de la nostra roba.

%%K t08-tit-2 sub
\VUsaut{5.0mm}
\VUpk{10.2pt}{40}{El Pastor.}
\VUsaut{3.4mm}

%%K t08-c1-05-1 p
5. --- El \VUgras{pastor} \textsuperscript{(10)} no sembla fer-los gaire cas. Només
es preocupa de la tempesta que passa per l'horitzó, i el \VUgras{sagal}
\textsuperscript{(13)}, que s'acosta la \VUgras{cantimplora}
\textsuperscript{(14)} als llavis, sembla encara més despreocupat.

%%K t08-c1-06-1 p
6. --- La calor és aclaparadora; perquè el \VUgras{gos} \textsuperscript{(15)} ve
també a refrescar-se; llepa l'aigua fresca del \VUgras{rierol}
\textsuperscript{(16)} i espanta una pobra \VUgras{granota}
\textsuperscript{(17)} que estava arraulida a l'herba de la riba. Aquesta salta a
l'aigua clara on floreixen els \VUgras{nenúfars} \textsuperscript{(18)}.

%%K t08-c1-07-1 p
7. --- Una \VUgras{cuereta} \textsuperscript{(19)} es banya vora la petita
\VUgras{font} \textsuperscript{(20)} que brolla entre dues roques i s'amaga sota la
\VUgras{brossa espinosa} \textsuperscript{(21)} i els \VUgras{matolls d'arç blanc}
\textsuperscript{(22)}.

%%K t08-tit-3 sub
\VUsaut{5.0mm}
\VUpk{10.2pt}{40}{La Collita.}
\VUsaut{3.4mm}

%%K t08-c1-08-1 p
8. --- Per als nens del camp, la sega del blat és una època molt divertida. Fan
alegres \VUgras{tombarelles} \textsuperscript{(24)} damunt els \VUgras{munts de
palla} \textsuperscript{(23)}, ajuden els \VUgras{segadors}
\textsuperscript{(26)} i, mentre aquests seguen el \VUgras{camp de blat}
\textsuperscript{(27)} amb les seves \VUgras{dalles} \textsuperscript{(25)}, ben
esmolades a la \VUgras{pedra d'esmolar} \textsuperscript{(30)}, ells recullen el
blat i el lliguen en \VUgras{garbes} \textsuperscript{(31)} o en \VUgras{feixos}
\textsuperscript{(32)}.

%%K t08-c1-09-1 p
9. --- De vegades es permet als grans de tallar amb una \VUgras{falç}
\textsuperscript{(12)} les \VUgras{tiges daurades} \textsuperscript{(28)} que
porten les pesades \VUgras{espigues} \textsuperscript{(29)} plenes de gra; i els
petits, mentre guarden el \VUgras{bestiar} \textsuperscript{(33)} que pastura al
\VUgras{prat} \textsuperscript{(34)} a l'ombra del gran \VUgras{til·ler}
\textsuperscript{(35)}, cacen les belles \VUgras{papallones} amb la seva
\VUgras{xarxa} \textsuperscript{(36)}.

%%K t08-c1-09-2 p
Alguns fins i tot es pugen al \VUgras{carro} \textsuperscript{(37)} per col·locar
les garbes que els passen amb una \VUgras{forca} \textsuperscript{(38)}.

%%K t08-c1-10-1 p
10. --- Per tota la \VUgras{plana} \textsuperscript{(39)}, on alcen el vol les
\VUgras{aloses} \textsuperscript{(41)}, hi regna l'animació. Tothom està ocupat amb
la sega. Però, mentre els pobres continuen usant les eines antigues ---
\VUgras{dalles, falçs} i \VUgras{batolls} \textsuperscript{(40)} ---, els
propietaris rics fan segar el seu blat o el seu \VUgras{sègol}
\textsuperscript{(43)} amb una \VUgras{segadora mecànica} \textsuperscript{(42)}.
Després, sense necessitat de portar les garbes al graner, s'hi fan allà mateix
grans \VUgras{garberes} \textsuperscript{(44)}.

%%K t08-c1-11-1 p
11. --- Aleshores es porta una \VUgras{batedora} \textsuperscript{(47)} que una
\VUgras{locomòbil} \textsuperscript{(45)} mou per mitjà d'una \VUgras{corretja de
transmissió} \textsuperscript{(46)}, i així el gra queda separat de la
\VUgras{palla} sense sortir del camp on va ser sembrat.

%%K t08-tit-4 sub
\VUsaut{5.0mm}
\VUpk{10.2pt}{40}{La Tempesta.}
\VUsaut{3.4mm}

%%K t08-c1-12-1 p
12. --- Sovint esclaten violentes \VUgras{tempestes} \textsuperscript{(53)} a
l'època de la sega. Primer se senten de lluny els trons del \VUgras{tro};
s'aixeca un fort \VUgras{vent de ponent} \textsuperscript{(54)} que doblega amb
gemecs els arbres més gruixuts del \VUgras{bosc} \textsuperscript{(49)}. Negres
\VUgras{nuvolots} \textsuperscript{(56)} assalten el cel; els travessen els
ziga-zagues encegadors del \VUgras{llamp} \textsuperscript{(55)}. La
\VUgras{pluja} i de vegades la \VUgras{calamarsa} comencen a caure, aixafant la
collita a punt de ser segada.

%%K t08-c1-13-1 p
13. --- Sovint el llamp incendia algun edifici. Així, l'any passat la teulada del
\VUgras{molí de vent} \textsuperscript{(50)} va quedar reduïda a cendres i dues de
les seves \VUgras{aspes} \textsuperscript{(51)} van quedar destrossades.

%%K t08-c1-14-1 p
14. --- Al cap d'unes hores torna la calma. Un brillant \VUgras{arc de Sant
Martí} \textsuperscript{(52)} apareix a l'horitzó i la natura reprèn la seva vida,
interrompuda un instant. Els pagesos, que s'havien refugiat a les
\VUgras{cabanes} \textsuperscript{(48)}, tornen a la feina i s'esforcen valentment
a reparar els danys que acaben de sofrir.

%%K t08-tit-5 sub
\VUsaut{6.0mm}
\VUcentre{13.2pt}{90}{\textit{Segona escena.}}
\VUsaut{3.6mm}

%%K t08-tit-6 sub
\VUpk{10.2pt}{40}{La Caça. --- La Verema.}
\VUsaut{1.4mm}
\VUpk{10.2pt}{40}{La Pesca.}
\VUsaut{4.0mm}

%%K t08-c2-01-1 p
1. --- Cap a finals d'\VUgras{agost} o a mitjan \VUgras{setembre}, segons les
regions, s'obre la veda. Amb prou feines han fet sortir els primers raigs del sol
les \VUgras{sargantanes} \textsuperscript{(2)} del seu forat, quan el
\VUgras{caçador} \textsuperscript{(5)} es posa en camí amb els seus
\VUgras{gossos} \textsuperscript{(4)}.

%%K t08-c2-02-1 p
2. --- S'ha posat el seu sòlid \VUgras{vestit de caça} \textsuperscript{(9)} de
drap gruixut i unes \VUgras{polaines} de cuir amb què podrà caminar per l'herba
molla on s'amaguen els \VUgras{gripaus} \textsuperscript{(1)}; ha agafat la seva
\VUgras{escopeta} \textsuperscript{(6)} i el seu \VUgras{sarró}
\textsuperscript{(10)}, i ja se n'ha anat.

%%K t08-c2-03-1 p
3. --- L'ardor de la persecució el porta a una vinya on la verema és en ple
apogeu. Els fills del vinyater, que solen guardar les cabres al camí, avui les
deixen pasturar al seu grat i, després de lligar a una estaca un vell
\VUgras{boc} \textsuperscript{(3)} que no deixaria d'aprofitar la seva llibertat
per escapar-se, es donen també la seva part de plaer. Un agafa un raïm d'una
\VUgras{portadora} \textsuperscript{(12)} i es diverteix fent córrer el
\VUgras{suc} \textsuperscript{(14)} en un \VUgras{got} \textsuperscript{(13)} per
fer most (vi sense fermentar). L'altre es corona amb una \VUgras{garlanda de
fulles} \textsuperscript{(15)} que acaba de trenar. Vora d'ells, uns
\VUgras{pardals} \textsuperscript{(16)} descarats arriben fins al mateix cup a
picotejar el raïm.

%%K t08-c2-04-1 p
4. --- Mentre els nens es diverteixen, els homes treballen. Els
\VUgras{veremadors} \textsuperscript{(18)} porten el raïm al celler en
\VUgras{cistells} \textsuperscript{(17)}; un \VUgras{boter}
\textsuperscript{(19)} arregla les \VUgras{bótes} \textsuperscript{(20)}, comprova
si les \VUgras{dogues} \textsuperscript{(22)} i els \VUgras{fons}
\textsuperscript{(21)} són en bon estat i substitueix els \VUgras{cèrcols}
\textsuperscript{(23)} trencats. També va ser ell qui va fer les grans
\VUgras{tines} \textsuperscript{(25)} en què s'aboca el \VUgras{raïm}
\textsuperscript{(26)} perquè fermenti.

%%K t08-c2-05-1 p
5. --- Quan la fermentació ha acabat, es treu el vi i es posa la resta al
\VUgras{cup} \textsuperscript{(28)} i, estrenyent a poc a poc el gran
\VUgras{cargol} \textsuperscript{(29)}, s'espremeix el suc, que corre a una
\VUgras{tineta} \textsuperscript{(32)}, i s'obté encara més \VUgras{vi}
\textsuperscript{(31)}.

%%K t08-tit-8 sub
\VUsaut{6.0mm}
\VUpk{10.2pt}{40}{Cervesa i Sidra.}
\VUsaut{3.4mm}

%%K t08-c2-06-1 p
6. --- Per la paret del \VUgras{celler} \textsuperscript{(27)} s'enfila una branca
de \VUgras{llúpol} \textsuperscript{(30)}. Aquí no és més que una planta d'adorn,
però en alguns països, on la vinya és molt rara, el llúpol es conrea per fabricar
la \VUgras{cervesa}.

%%K t08-c2-07-1 p
7. --- La petita \VUgras{pomera} \textsuperscript{(33)} és carregada de pomes que,
quan siguin madures, donaran una \VUgras{sidra} excel·lent. A la seva
\VUgras{gàbia} \textsuperscript{(34)}, el \VUgras{canari} \textsuperscript{(35)} i
la \VUgras{cadernera} \textsuperscript{(36)} miren amb enveja els pardals que
s'esbargeixen en llibertat.

%%K t08-c2-08-1 p
8. --- Fins on abasta la vista, al vessant dels turons, s'alineen els
\VUgras{rodrigons} \textsuperscript{(40)} que sostenen els magnífics \VUgras{ceps}
\textsuperscript{(39)}, carregats de pesants \VUgras{raïms}.

%%K t08-c2-08-2 p
Durant tot l'any, el \VUgras{vinyater} \textsuperscript{(42)} ha treballat per
tenir cura de la seva \VUgras{vinya} \textsuperscript{(38)}, i ara es veu
recompensat del seu esforç amb una bella collita. Les \VUgras{veremadores}
\textsuperscript{(41)} omplen les tines posades al carro que tiren les
\VUgras{mules} \textsuperscript{(43)}, i tothom està alegre.

%%K t08-tit-9 sub
\VUsaut{5.0mm}
\VUpk{10.2pt}{40}{La Pesca.}
\VUsaut{3.4mm}

%%K t08-c2-09-1 p
9. --- El mateix passa amb els \VUgras{pescadors de canya}
\textsuperscript{(44)}, que des del matí han vingut a instal·lar-se a la vora de
l'aigua amb les seves \VUgras{canyes de pescar} \textsuperscript{(45)}.

%%K t08-c2-09-2 p
El \VUgras{peix} \textsuperscript{(47)} pica a l'ham tan bon punt tiren el seu
\VUgras{fil} \textsuperscript{(46)} a l'aigua, i amb prou feines els dóna temps de
renovar l'\VUgras{esquer} quan una nova víctima cueja a la punta del fil.

%%K t08-c2-09-3 p
Amb tot, el \VUgras{pescador de xarxa} \textsuperscript{(48)}, que des de la seva
\VUgras{barca} \textsuperscript{(49)} tira l'esparver al mig del \VUgras{riu}
\textsuperscript{(50)}, n'agafa molts més.

%%K t08-c2-10-1 p
10. --- La tardor és l'estació dels bons passeigs: a peu, a \VUgras{cavall}
\textsuperscript{(52)}, en \VUgras{bicicleta} \textsuperscript{(53)}, en
\VUgras{automòbil} \textsuperscript{(54)}. Els \VUgras{camins}
\textsuperscript{(51)} són nets i s'aprofiten els últims dies bons, perquè ja les
\VUgras{orenetes} \textsuperscript{(56)} es reuneixen a les teulades per volar a
tota ala cap a països més càlids. Les fulles de la vella \VUgras{noguera}
\textsuperscript{(57)} groguegen, les \VUgras{nous} cauen i els alts
\VUgras{arbres} agrupats en \VUgras{bosquet} \textsuperscript{(58)} damunt
l'\VUgras{altura} \textsuperscript{(59)} aviat es quedaran sense fulles.

%%K t08-c2-11-1 p
11. --- El \VUgras{sol} \textsuperscript{(60)} surt més tard, els dies minven,
comença a sentir-se un fred prou viu al \VUgras{matí}, i ja bandades de
\VUgras{becades} \textsuperscript{(61)} deixen els nostres climes inhospitalaris.
\VUsaut{6.0mm}
\VUfilet{22mm}
